\chapter*{Appendix R: Stratified Semantic Manifolds and Hypergraph Fibrations}
\addcontentsline{toc}{chapter}{Appendix R: Stratified Semantic Manifolds and Hypergraph Fibrations}

\section{Overview}

This appendix develops the topological and differential-geometric foundations of the semantic manifolds that support the RSVP field theory and the Polyxan hypergraph apparatus. The construction unifies smooth stratified curvature manifolds, quine-stable hypergraph bundles, and the canonical embeddings that link them.

Whereas the body of the monograph treats these structures in operational form, the present appendix provides a unified mathematical account of their topology, stratification, fibration, and compatibility conditions. This is the geometric substrate upon which modal semantics, categorical duality, and verification are built.

\section{Stratified Semantic Manifolds}

\subsection{Semantic Manifold}

A semantic manifold is a quadruple
\[
\mathcal{M} = (M,g,\mathcal{S},\mathcal{R})
\]
where:
\begin{itemize}
\item $M$ is a connected, second-countable topological manifold,
\item $g$ is a piecewise-smooth RSVP metric,
\item $\mathcal{S} = \{M_k\}_{k=0}^n$ is a Whitney stratification,
\item $\mathcal{R} : M \to \mathbb{Z}_{\ge 0}$ assigns an integer curvature stratum index.
\end{itemize}

Each stratum $M_k$ is a smooth manifold of codimension $k$. The curvature index $\mathcal{R}$ is constant on strata and strictly increases with codimension.

\subsection{Curvature Layering}

The RSVP curvature scalar is defined on each stratum and extended continuously to the stratification by:
\[
\mathcal{R}(x) = k
\quad\Longleftrightarrow\quad
x \in M_k.
\]
This definition ensures monotonic descent under harmonic flow and stratified invariance under curvature-matched embeddings.

\section{Polyxan Hypergraph Fibrations}

\subsection{Incidence Fibration}

Let $\mathcal{G}$ be a finite typed Polyxan hypergraph. Its incidence set
\[
\mathrm{Inc}(\mathcal{G})
\]
forms the total space of a discrete fibration
\[
\pi : \mathrm{Inc}(\mathcal{G}) \to |\mathcal{G}|
\]
where $|\mathcal{G}|$ is the topological realisation of the hypergraph.

\subsection{Quine-Stable Fibres}

If $\mathcal{G}$ is quine-stable, the fibres of $\pi$ are saturated under EHR rewriting: each fibre contains an entire confluence class, and the projection respects rewrite equivalence. This structure is essential for the embedding compatibility discussed below.

\section{Canonical Embeddings}

\subsection{Definition}

An embedding of a hypergraph into the semantic manifold is a map
\[
\iota : \mathrm{Inc}(\mathcal{G}) \hookrightarrow M
\]
satisfying the curvature-matching and entropy-alignment equations:
\[
\mathcal{R}\circ \iota = \mathcal{H}_0[\mathcal{G}],
\qquad
S\circ \iota = \mathcal{H}_{\mathrm{tag}}[\mathcal{G}].
\]

\subsection{Stratum Compatibility}

The embedding is stratified if:
\[
\iota(M_k) \subseteq M_k
\]
for all $k$. This ensures that the discrete curvature of $\mathcal{G}$ exactly matches the geometric curvature strata of the manifold.

\subsection{Harmonicity}

An embedding is harmonic when its Dirichlet energy vanishes on each stratum, equivalently:
\[
\Delta_g \iota = 0
\quad\text{on } M_k.
\]
This condition ensures uniqueness up to stratum-preserving isometries.

\section{Semantic Bundles}

\subsection{Universal Hypergraph Bundle}

Over each stratum $M_k$ define a bundle:
\[
\mathcal{U}_k = M_k \times \mathcal{C}_k
\]
where $\mathcal{C}_k$ is the class of hypergraphs of curvature index $k$.

Gluing the bundles across the stratification yields a global semantic bundle:
\[
\mathcal{U} = \bigsqcup_{k=0}^n \mathcal{U}_k.
\]

\subsection{Sections}

A semantic world is a section of the universal bundle satisfying:
\[
\sigma(x) = (\; x,\; \mathcal{G}_x \;)
\]
with $\mathcal{G}_x$ required to vary continuously across strata in the sense of rewrite fibration equivalence.

\section{Consistency Conditions}

\subsection{Stratified Coherence}

A section is coherent when the following holds for any path $\gamma$ crossing strata:
\[
\mathcal{G}_{\gamma(0)} \quad\rightsquigarrow\quad \mathcal{G}_{\gamma(1)}
\]
via a finite sequence of rewrite and quine-stabilisation steps, ensuring that semantic transitions preserve hypergraph type.

\subsection{Curvature–Rewrite Compatibility}

For any $x,y \in M$,
\[
x \to y \text{ by curvature descent}
\quad\Longrightarrow\quad
\mathcal{G}_x \rightsquigarrow \mathcal{G}_y \text{ by rewrite descent}.
\]

This captures the geometric–algebraic duality at the heart of the monograph.

\section{Topos of Semantic Structures}

\subsection{Category of Worlds}

Define the category:
\[
\mathbf{Sem} = \mathbf{Sh}(\mathcal{U})
\]
whose objects are sheaves on the semantic bundle and morphisms are sheaf morphisms.

\subsection{Internal Logic}

The internal logic of $\mathbf{Sem}$ is intuitionistic and supports the modal operators introduced later. Truth values correspond to geometric invariants across strata.

\subsection{Geometric Morphisms}

The RSVP and Polyxan modalities each induce a geometric morphism:
\[
\Box_{\mathrm{RSVP}},\quad \Box_{\mathrm{Polyx}} : \mathbf{Sem} \to \mathbf{Sem}
\]
defined by inverse image functors along curvature and rewrite fibrations, respectively.

\section{Connections to the Main Text}

The structures introduced here provide:
\begin{itemize}
\item the manifold topology used in curvature stratification (Ch.~13),
\item the hypergraph fibrations used in semantic galaxy formation (Ch.~30),
\item the harmonic embeddings used in the synthetic duality (Ch.~62),
\item the bundle structure used in the reconciliation algebra (Ch.~79),
\item the geometric basis for modal semantics (Appendix S).
\end{itemize}

\section{Summary}

This appendix presented:
\begin{itemize}
\item stratified semantic manifolds and curvature layering,
\item Polyxan hypergraph fibrations and quine-stable fibres,
\item canonical stratified embeddings with harmonicity conditions,
\item the universal semantic bundle and its internal topos,
\item the foundational structures needed to interpret modal necessity geometrically.
\end{itemize}

These constructions form the geometric foundation upon which subsequent modal, categorical, and epistemic results rely.

